\SourceLeaf{226}%
\bigskip\Needspace{9\baselineskip}
\pdfbookmark[1]{Pallas Olbersiana. 1802 August}{pallas180208}
\begin{center}
PALLAS OLBERSIANA.\par
\medskip\hrule\medskip
Monatliche Correspondenz. Band VI. S.\ 187—196. 1802 August.\par
\medskip\hrule
\end{center}\medskip
—\quad —\quad —\quad Auch aus \textit{Vilna}, \textit{Prag}, \textit{Kremsmünster} und \textit{Padua} haben wir die Bestätigung erhalten, dass \textit{Pallas} dort nicht beobachtet wurde. Nach einer Nachricht von Dr.\ \textsc{Gauß} hatte Dr.\ \textsc{Maskelyne} die Beobachtung dieses Planeten am 20. Juli noch nicht aufgegeben. Er glaubt, \textit{Pallas} am 18. Juli noch als Stern 11. Größenklasse beobachtet zu haben, war sich aber nicht ganz sicher, ob es der Planet oder ein Fixstern war. Das wird die Rechnung oder die nächste Beobachtung zeigen.

Dr.\ \textsc{Gauß}’ dritte Bahnelemente für \textit{Pallas} stimmen weiterhin gut mit den beobachteten Deklinationen überein. Nur die Rektaszensionen weichen jetzt um einige Sekunden mehr ab. Dr.\ \textsc{Gauß} verglich die Beobachtungen, die \textsc{Cesaris} in Mailand an einem ausgezeichneten achtfüßigen Mauerquadranten von \textsc{Ramsden} angestellt hatte und die wir im vorigen Heft auf S.\ 73 mitgeteilt haben. Er erhielt die folgende Übereinstimmung.
\begin{center}
\textit{Vergleich der Mailänder Meridianbeobachtungen der Pallas}\par
\textit{mit den dritten Elementen von }\textsc{Gauß}.\par\medskip
{\small\setlength{\tabcolsep}{4pt}
\begin{tabular}{r r|r|r|r|r|r}
\multicolumn{2}{c|}{1802}&\shortstack{Mittlere Zeit\\in Mailand}&\multicolumn{2}{c|}{Berechnete}&\multicolumn{2}{c}{\shortstack{Unterschied\\R—B}}\\
\multicolumn{2}{c|}{}&&\shortstack{Rektaszension}&\shortstack{Nördliche\\Deklination}&in R.A.&in Dekl.\\\hline
Mai&4&$9^{\mathrm{u}}\ 15^{\mathrm{m}}\ 43^{\mathrm{s}}1$&$180^0\ 56'\ 15''7$&$20^0\quad2'\ 44''1$&$+15''7$&$-5''9$\\
&5&$9\ 11\ 45.3$&$180\ 55\ 56.8$&$20\quad8\ 38.1$&$+19.8$&$+5.1$\\
&6&$9\quad7\ 49.3$&$180\ 56\quad0.3$&$20\ 14\ 10.4$&$+16.3$&$+6.4$\\
&7&$9\quad3\ 55.6$&$180\ 56\ 24.8$&$20\ 19\ 22.1$&$+10.8$&$+2.1$\\
&8&$9\quad0\quad3.2$&$180\ 57\ 13.3$&$20\ 24\ 12.4$&$+8.3$&$+0.4$\\
&9&$8\ 56\ 12.1$&$180\ 58\ 21.9$&$20\ 28\ 43.1$&$+6.9$&$+0.1$\\
&10&$8\ 52\ 22.4$&$180\ 59\ 52.8$&$20\ 32\ 53.0$&$+7.8$&$-2.0$\\
&11&$8\ 48\ 34.3$&$181\quad1\ 45.2$&$20\ 36\ 43.4$&$+9.2$&$-3.6$\\
&12&$8\ 44\ 47.0$&$181\quad3\ 58.8$&$20\ 40\ 14{,}6$&$+7.8$&$-1.4$\\
&17&$8\ 26\ 12.9$&$181\ 20\ 17.1$&$20\ 53\ 18.0$&$+11.1$&$-4.0$\\
&18&$8\ 22\ 33.1$&$181\ 24\ 33.5$&$20\ 55\quad3.6$&$+15.5$&$+0.6$\\
&19&$8\ 18\ 55.8$&$181\ 29\quad9.4$&$20\ 56\ 30.8$&$+16.4$&$-0.2$\\
&20&$8\ 15\ 19.6$&$181\ 34\quad4.8$&$20\ 57\ 43.0$&$+14.8$&$-2.0$\\
&21&$8\ 11\ 44.6$&$181\ 39\ 18.6$&$20\ 58\ 39.5$&$+10.6$&$-1.5$\\
&22&$8\quad8\ 11.1$&$181\ 44\ 52.9$&$20\ 59\ 20.6$&$+6.9$&$-1.4$
\end{tabular}}
\end{center}

\Needspace{22\baselineskip}
Auch die Beobachtungen von Dr.\ \textsc{Olbers} zeigen dieselbe Übereinstimmung. Dass sie auch untereinander so gut übereinstimmen, zeigt zugleich den besonderen Wert dieser Beobachtungen, obwohl sie nur am \textit{Kreismikrometer} angestellt wurden. Wer sich die Mühe macht, sie ebenso genau zu untersuchen wie wir, wird feststellen, dass sie denen an Äquatorialsektoren in nichts nachstehen, ihnen ebenbürtig und mitunter sogar vorzuziehen sind. Das spricht freilich mehr für die Geschicklichkeit des ausgezeichneten Beobachters als für die Güte des Instruments, wie unsere Leser aus dem folgenden Vergleich erkennen werden.
\newcommand{\PallasSymbol}{\text{\mdseries\AstroSymbols ⚴}}
\SourceLeaf{227}%
\Needspace{15\baselineskip}
\begin{center}
\textit{Vergleich der Pallasbeobachtungen von Dr.\ }\textsc{Olbers}\par
\textit{mit den dritten Elementen von }\textsc{Gauß}.\par\medskip
{\small\setlength{\tabcolsep}{4pt}
\begin{tabular}{r r|r|r|r|r|r}
\multicolumn{2}{c|}{1802}&\shortstack{Mittlere Zeit\\in Bremen}&\multicolumn{2}{c|}{Berechnete}&\multicolumn{2}{c}{\shortstack{Unterschied\\R—B}}\\
\multicolumn{2}{c|}{}&&\shortstack{Rektaszension\\der $\PallasSymbol$}&\shortstack{Rektaszension\\der $\PallasSymbol$}&in R.A.&in Dekl.\\\hline
Juni&19&$11^{\mathrm{u}}\ 13^{\mathrm{m}}\ 55^{\mathrm{s}}$&$186^0\ 16'\ 19''$&$19^0\ 54'\ 17''$&$-6''$&$+17''$\\
&20&$11\quad7\ 51$&$186\ 29\ 12$&$19\ 49\ 39$&$-22$&$-1$\\
&21&$11\quad8\ 54$&$186\ 42\ 39$&$19\ 44\ 55$&$-20$&$-3$\\
&26&$11\ 18\ 34$&$187\ 52\quad3$&$19\ 19\ 31$&$-16$&$+8$\\
Juli&4&$11\ 28\ 14$&$189\ 52\ 12$&$18\ 33\ 39$&$-26$&$-5$\\
&8&$11\quad5\ 16$&$190\ 55\ 44$&$18\quad8\ 50$&$-25$&$-8$\\
&8&$11\ 13\ 20$&$190\ 55\ 50$&$18\quad8\ 48$&$-23$&$-10$\\
&9&$11\ 53\ 47$&$191\ 12\ 36$&$18\quad2\ 35$&$-28$&$+1$
\end{tabular}}
\end{center}

Aus all diesen Vergleichen und der ausgezeichneten Übereinstimmung der Beobachtungen mit der Ellipse von \textsc{Gauß} (die dieser geschickte Rechner der Wahrheit noch näher bringen wird) dürfen wir zuversichtlich hoffen, dass uns dieser neue Planet selbst dann erhalten bleibt, wenn er wegen seiner Lichtschwäche 1803 und 1804 unmöglich zu finden sein sollte. Doch schon jetzt haben wir begründete Hoffnung, dass sich \textit{Pallas} im kommenden Jahr unseren geschärften Blicken nicht entziehen wird. Dr.\ \textsc{Olbers}, der diesen Planeten zuletzt am 9. Juli gesehen hatte, ist selbst dieser Meinung und schreibt dazu. „So schwer der Planet auch jetzt zu sehen ist, bin ich doch überzeugt, dass Sie ihn im nächsten Jahr an Ihrem Durchgangsinstrument sehen werden. Bei seiner Opposition ist er der Erde dann etwas näher, als er ihr am 9. Juli war. Zwar steht er weiter von der Sonne entfernt und wird von ihr schwächer beleuchtet, doch seine größere Höhe über dem Horizont und der größere Abstand von der Dämmerung gleichen das mehr als aus.“

\Needspace{24\baselineskip}
Auch Dr.\ \textsc{Gauß} ist sehr zuversichtlich, dass sich \textit{Pallas} im nächsten Jahr wiederfinden lässt. Für den 28. Jan.\ 1803 (etwa einen Tag vor der Opposition) ermittelt er ihren Abstand von der Erde zu $2.56$. Aus Dr.\ \textsc{Herschel}s eigener Messung ermittelt er für diesen Planeten einen etwas größeren Durchmesser. Am 22. April gibt Dr.\ \textsc{Herschel} nach einer recht guten Messung den scheinbaren Durchmesser mit $=0''17$ an. Daraus berechnet Dr.\ \textsc{Gauß} den wahren Durchmesser zu $26\frac{1}{2}$ deutschen Meilen (bei einem Abstand von der Erde $=1.562$). In seinem letzten Brief wundert er sich darüber, wie außerordentlich verschieden die Ergebnisse der Messungen dieser Durchmesser durch Dr.\ \textsc{Herschel} und Dr.\ \textsc{Schröter} sind, obwohl sie nach derselben Methode vorgenommen wurden. „Ich bin sehr neugierig“, schreibt er, „wie stark Dr.\ \textsc{Herschel} vergrößert hat. Bei 500-facher Vergrößerung würde ein scheinbarer Durchmesser von $0''17$ wohl noch keine Scheibe ergeben?“ Ich jedenfalls habe bei 300-facher Vergrößerung weder an diesem Planeten von \textsc{Olbers} noch an jenem von \textsc{Piazzi} eine Spur einer Scheibe wahrnehmen können. Auf meine Anfrage dazu antwortete Dr.\ \textsc{Olbers} wie folgt. „Auch ich habe selbst bei den stärksten Vergrößerungen meines \textsc{Dollond}-Achromaten nie einen Unterschied zwischen diesen Planeten und Fixsternen von gleicher scheinbarer Größe wahrnehmen können.“

Wir haben im vorigen Heft auf S.\ 85 [S.\ 219 dieses Bandes] eine Ephemeride für \textit{Pallas} von \textsc{Gauß} mitgeteilt. Hier folgt ihre Fortsetzung. Sie erscheint zwar zu spät, um die Beobachtungen noch zu erleichtern, wird aber doch dazu beitragen, dass Sterne, die man etwa mit dem Planeten verwechselt hat, nicht als \textit{Pallas} ausgegeben werden. [Die Ephemeride ist ebenso eingerichtet wie die auf S.\ 219 abgedruckte, enthält nur den Abstand der Pallas von der Sonne nicht und reicht bis zum 28. August 1802. Die bisher in dieser Ausgabe von \textsc{Gauß}’ Werken mitgeteilten Ephemeriden werden als Beispiele für die von ihm berechneten genügen. Die übrigen können aus Platzgründen unberücksichtigt bleiben.]
